\clearpage
\section{Resumen ejecutivo}

\subsection{Presentación del proyecto}

El proyecto desarrolla en doce meses un MVP portátil para seguimiento metabólico complementario de personas con diabetes tipo 2 o prediabetes. Integra un parche sensor descartable, un módulo lector reutilizable, una pulsera reutilizable y una aplicación móvil. El sistema busca reunir biomarcadores en sudor y señales fisiológicas, presentar historial e incertidumbre e investigar factores probablemente asociados a las variaciones observadas. No reemplaza un CGM, no establece causalidad clínica y no emite recomendaciones automáticas de tratamiento.

El equipo base es de cuatro personas full-time ---dos perfiles de ingeniería y dos técnicos/administrativos---, equivalentes a 7.680 horas brutas de capacidad anual. El prototipo integrado del mes 7 es un hito intermedio; el resultado final es un MVP aceptado funcional y técnicamente en el mes 12. La validación clínica y la certificación comercial no están incluidas.

\subsection{Mercado, costos y precios}

La referencia de población es de aproximadamente cuatro millones de adultos con diabetes tipo 2 en Argentina, pero no se la considera demanda capturable. La entrada al mercado se modela mediante médicos y clínicas aliados. En el escenario base se prevén 18/35/56 médicos y 7/15/24 clínicas en los años 2 a 4, alcanzados desde 90/175/280 prospectos médicos y 35/75/120 clínicas calificadas con 20\% de conversión. Con 40, 50 y 60 pacientes alcanzables por aliado y tasas de adopción de 30\%, 32\% y 35\%, se obtienen 300, 800 y 1.680 kits vendidos. Los médicos son el canal principal; cada clínica recibe un demostrador y onboarding institucional desde el año 2. El precio de planificación es USD 400 por kit y USD 25 por parche de recambio; la retención mensual supuesta es 95\%.

El presupuesto del año 1 es USD 228.303,06. Incluye doce meses de dotación, infraestructura, servicios generales, laboratorio/equipamiento, prototipos, ensayos, consultas especializadas, QA, herramientas, preserie y una reserva técnica trazable. En explotación, el costo directo de fabricación es USD 102,46 por kit y la logística/importación USD 69,06. En la preserie, el costo directo baja a USD 87,20 y la logística/importación a USD 64,03 porque la terminación, programación, prueba y preparación del parche ya están cubiertas por los cuatro FTE. En ambos casos se aplica una carga académica simplificada del 45\% sobre CIF. No se duplican salarios con horas valorizadas ni se utiliza un coeficiente general de fabricación.

\subsection{Diseño técnico y aspectos a cubrir}

El MVP debe verificar adquisición, calibración, memoria, sincronización BLE, recuperación de datos, autonomía, montaje, recambio, robustez, repetibilidad y trazabilidad de lotes. Como base de diseño, cada dispositivo reserva al menos 128 KiB utilizables para un buffer de 24 horas y apunta a siete días de autonomía, sujeto a perfil de consumo y pruebas de descarga. Compras, importación, despacho, recepción, inventario, QA, soporte y documentación son parte del alcance del proyecto.

\subsection{FODA y riesgos}

Las fortalezas internas son la dotación full-time con roles definidos, la capacidad multidisciplinaria de integración, la gestión trazable de requisitos/hitos/pruebas y el modelo operativo-financiero conectado. Las principales debilidades son la experiencia práctica e infraestructura especializadas todavía limitadas y la falta de desempeño propio demostrado. Los riesgos dominantes son integración/calibración, disponibilidad de datos, abastecimiento, plazo, adopción, costos y definición regulatoria. Se revisan mensualmente con una rúbrica común de severidad y ocurrencia y con evidencia de mitigación.

\subsection{Flujo de fondos, financiamiento y viabilidad}

El año 1 es exclusivamente de desarrollo e inversión y no registra ventas. El lanzamiento se inicia en el año 2, condicionado al MVP y a los requisitos regulatorios aplicables. El capital mínimo para cubrir el mayor déficit acumulado es USD 228.303,06. Con los supuestos base, el VAN a una tasa anual del 20\% es USD 103.250,74 y la TIR anual 36,83\%. No se aumentaron las ventas para obtener esos indicadores: el cambio proviene del costeo trazable. El escenario conservador conserva un VAN negativo de USD 223.190,04 y la prueba con retención mensual del 85\% da VAN negativo de USD 18.477,62; por eso la viabilidad depende de validar el embudo, el margen, los recambios y la importación antes de escalar.
